% =====================================================================
% Commandes.tex
% Custom math shortcuts, vectors, integrals, and physics-specific macros.
% Part of the definitions-latex package.
%
% Required packages: amsmath, amssymb, bm, ifthen, hyperref, mhchem, bbold
% =====================================================================

% ---------------------------------------------------------------------
% Equation tagging
% ---------------------------------------------------------------------
\newcommand{\numberthis}{\addtocounter{equation}{1}\tag{\theequation}}

% ---------------------------------------------------------------------
% Hyperlinks
% ---------------------------------------------------------------------
\newcommand{\mail}[1]{\href{mailto:#1}{#1}}
\newcommand{\ftplink}[1]{\href{ftp://#1}{#1}}

% ---------------------------------------------------------------------
% Differentials and partial derivatives
% ---------------------------------------------------------------------
\newcommand{\dif}{\mathrm{d}}
\newcommand{\fracdif}[2]{\frac{\dif#1}{\dif#2}}
\newcommand{\dift}{\frac{\partial}{\partial{t}}}
\newcommand{\partialFrac}[2]{\frac{\partial#1}{\partial#2}}
\newcommand{\fp}[2]{\frac{\partial#1}{\partial#2}}
\newcommand{\parfrac}[2]{\frac{\partial#1}{\partial#2}}
\newcommand{\parfracu}[1]{\parfrac{}{#1}}
\newcommand{\parfracR}{\parfracu{r}}

% ---------------------------------------------------------------------
% Vectors and unit vectors
% ---------------------------------------------------------------------
\newcommand{\vo}{\bm{0}}
\newcommand{\ex}{\bm{e}_x}
\newcommand{\ey}{\bm{e}_y}
\newcommand{\ez}{\bm{e}_z}
\newcommand{\er}{\bm{e}_r}
\newcommand{\ep}{\bm{e}_\varphi}
\newcommand{\et}{\bm{e}_{\theta}}
\newcommand{\x}{\bm{x}}
\newcommand{\X}{\bm{X}}
\newcommand{\vg}{\bm{v}}
\newcommand{\ug}{\bm{u}}
\newcommand{\grad}{\bm{\nabla}}

% Column vector of two entries
\newcommand{\vcold}[2]{\begin{bmatrix} #1 \\ #2 \end{bmatrix}}

% ---------------------------------------------------------------------
% Brackets, norms, grouping
% ---------------------------------------------------------------------
\newcommand{\Abs}[1]{\left|#1\right|}
\newcommand{\Absfrac}[2]{\Abs{\frac{#1}{#2}}}
\newcommand{\Absfracu}[1]{\Absfrac{\Delta#1}{#1}}
\newcommand{\nor}[1]{\left\Vert#1\right\Vert}
\newcommand{\Par}[1]{\left(#1\right)}
\newcommand{\Cro}[1]{{\left[#1\right]}}
\newcommand{\mean}[1]{\langle#1\rangle}
\newcommand{\eval}[1]{\left.#1\right|}

% ---------------------------------------------------------------------
% Integrals
% ---------------------------------------------------------------------
\newcommand{\intt}{\int{}\!\!\!\int{}\!\!\!\int} % \! is intended; chktex 21
\newcommand{\diftt}{\dif{x}\dif{y}\dif{z}}

% ---------------------------------------------------------------------
% Statistics shorthands
% ---------------------------------------------------------------------
\newcommand{\median}[1]{#1^{(\lceil{N}/2\rceil)}}
\newcommand{\iid}{\overset{\mathrm{iid}}{\sim}}

% ---------------------------------------------------------------------
% Common shorthands
% ---------------------------------------------------------------------
\newcommand{\oover}[1]{\frac{1}{#1}}
\newcommand{\xsover}[1]{\frac{x^*}{#1}}
\newcommand{\kro}{\delta_{ab}}
\newcommand{\Reyn}{\mathit{Re}}
\newcommand{\comma}{,}                       % useful inside siunitx args
\newcommand{\tensor}[1]{\stackrel{\leftrightarrow}{#1}}

% ---------------------------------------------------------------------
% Identity matrix (requires the bbold package)
% ---------------------------------------------------------------------
\newcommand{\iden}{\text{\usefont{U}{bbold}{m}{n}1}}
\MakeRobust{\iden}

% =====================================================================
% Plasma-physics specific shorthands
% =====================================================================
\newcommand{\Zeff}{Z_{\mathit{eff}}}
\newcommand{\OH}{\mathit{OH}}
\newcommand{\kb}{k_b}
\newcommand{\Beddy}{B_\mathit{eddy}}
\newcommand{\prad}{P_{\mathit{rad}}}
\newcommand{\refMath}{\mathit{ref}}
\newcommand{\Iref}{I_\refMath}
\newcommand{\vloop}{V_\mathit{loop}}
\newcommand{\wall}{\mathit{wall}}
\newcommand{\ode}{\mathit{ODE45}}

% Half-step temperature
\newcommand{\Te}{T + \Delta{T} / 2}

% ITER open-loop scenario
\newcommand{\OL}[1]{\mathit{ITER#1}_{OL}}

% Initial densities (uses mhchem \ce)
\newcommand{\initialdensity}[1]{n_{\ce{#10}}}
\newcommand{\correctionDeuteriumDensity}{C_{\initialdensity{D}}}

% Yield notation Y^{species A}_{species B}, with all 3x3 D/C/O combinations
\newcommand{\YAsB}[2]{Y^{\ce{#1}}_{\ce{#2}}}
\newcommand{\YDsD}{\YAsB{D}{D}}
\newcommand{\YDsC}{\YAsB{D}{C}}
\newcommand{\YDsO}{\YAsB{D}{O}}
\newcommand{\YCsD}{\YAsB{C}{D}}
\newcommand{\YCsC}{\YAsB{C}{C}}
\newcommand{\YCsO}{\YAsB{C}{O}}
\newcommand{\YOsD}{\YAsB{O}{D}}
\newcommand{\YOsC}{\YAsB{O}{C}}
\newcommand{\YOsO}{\YAsB{O}{O}}

% ---------------------------------------------------------------------
% Goldstone confinement-time scaling law
% Args: 1 = density subscript, 2..9 = exponents on B_phi, I_p, n_e, a,
%                                       R_0, kappa, A, P_Omega.
% A zero exponent suppresses the corresponding factor entirely.
% ---------------------------------------------------------------------
\newcommand{\GoldstoneScalinglaw}[9]{%
{\ifthenelse{\equal{#2}{0}}{}{B_\phi^{#2}}}\,
{\ifthenelse{\equal{#3}{0}}{}{I_{p,\,\mathrm{MA}}^{#3}}}\,
{\ifthenelse{\equal{#4}{0}}{}{n_{e,#1}^{#4}}}\,
{\ifthenelse{\equal{#5}{0}}{}{a^{#5}}}\,
{\ifthenelse{\equal{#6}{0}}{}{R_0^{#6}}}\,
{\ifthenelse{\equal{#7}{0}}{}{\kappa^{#7}}}\,
{\ifthenelse{\equal{#8}{0}}{}{A^{#8}}}\,
{\ifthenelse{\equal{#9}{0}}{}{P_{\Omega,\,\mathrm{MW}}^{#9}}}%
}

% =====================================================================
% Text helpers
% =====================================================================

% C++ logo (kept once -- the original file defined it twice)
\newcommand{\CC}{C\nolinebreak[4]\hspace{-.05em}\raisebox{.4ex}{\tiny\normalfont\bfseries++}}

% Roman numerals: \rom{4} -> IV
\newcounter{romanNumeralCounter}

\newcommand*{\rom}[1]{%
    \setcounter{romanNumeralCounter}{#1}%
    \Roman{romanNumeralCounter}%
}

% =====================================================================
% Custom appendices: "Appendix I", "Appendix II", ...
% =====================================================================
\newcounter{myappendix}
\newcommand{\myappendix}[1]{%
  \refstepcounter{myappendix}%
  \section*{Appendix~\Roman{myappendix}: #1}%
}

% =====================================================================
% Section numbering depth (number down to \paragraph level)
% =====================================================================
\setcounter{secnumdepth}{4}
